% !TeX root = iclr2027_conference.tex

% Main-text Introduction. Citation keys are defined in iclr2027/iclr2027_conference.bib.

\section{Introduction}
\label{sec:introduction}
Large language models (LLMs) have achieved strong performance across a broad range of tasks \citep{openai2024gpt4,geminiteam2025gemini}, but capability does not guarantee factual reliability.  TruthfulQA shows that models can reproduce common human falsehoods \citep{lin2022truthfulqa}, while HaluEval documents hallucinated content across knowledge-intensive question answering, dialogue, and summarization \citep{li2023halueval}.  This problem becomes more consequential as LLMs are embedded in tool-augmented and agentic workflows, including web search and external-service interaction \citep{nakano2022webgpt,liu2023webglm,steinberger2026openclaw}.  Reliably detecting hallucinations is therefore essential before model outputs trigger real-world actions.

One particularly difficult case is a factually incorrect answer produced with high confidence even though the required knowledge is behaviorally accessible.  Existing detectors provide limited insight into such errors.  Uncertainty-based methods use likelihood, self-evaluation, or disagreement across samples \citep{kadavath2022mostlyknow,manakul2023selfcheckgpt,farquhar2024semanticentropy}, but systematic errors can be confident and repeatable.  Hidden-state and attention-based detectors identify representations correlated with incorrectness \citep{azaria2023internal,chen2024inside,du2024haloscope,chuang2024lookback,binkowski2025spectral}, but such correlations may reflect dataset artifacts \citep{janiak2025illusion}.  Evidence-grounded methods verify generated claims against retrieved sources \citep{min2023factscore}, but depend on external evidence and assess support rather than the cause of the model's decision.  None of these approaches reveals which input evidence drives a confident, repeatable mistake.

Prior work suggests that these errors can arise when statistical associations override prompt-supported reasoning \citet{mckenna2023sources, sun2025subsequence}. Most directly, \citet{hu2026reasoningpriors} formalize this phenomenon as inference misalignment in a latent key--task model and introduce ScientistQA to distinguish missing knowledge from failures to deploy accessible facts.  Their analysis establishes a theoretical account and demonstrates the phenomenon behaviorally, but does not resolve how a model's internal computation selects and routes the competing evidence in a particular prompt.  

In this work, we investigate this gap through targeted keyword perturbations.  We perturb one keyword at a time in the model's input representation and measure how its relative preference between candidate answers changes.  If perturbing a keyword reduces the model's preference for a wrong answer, that keyword likely contributed to the hallucination; if it instead weakens the correct answer, the keyword provides evidence the model was using appropriately.  The resulting pattern offers an interpretable signal for both hallucination detection and attribution of individual errors.  Experiments across three LLMs and multiple benchmarks show that this signal generalizes across tasks and transfers to new domains.

We also examine the operating regimes of four common families of hallucination detectors: uncertainty-based, representation-based, evidence-based, and keyword-perturbation methods.  Our results suggest that their effectiveness depends on the type of error, which may help explain why detectors do not always transfer reliably across benchmarks.  Based on the observed patterns, we organize hallucinations into four broad categories: missing knowledge, context-induced errors, self-consistent false knowledge, and unstable reasoning.  Because the prevalence of these categories varies across benchmarks, no single detection method should be expected to perform uniformly well in every setting.

\paragraph{Contributions.}
Our main contributions are summarized as follows:
\begin{enumerate}
    \item We investigate why models hallucinate even when the relevant knowledge has been acquired.  We show that model predictions can depend on particular keywords that are bound to the question entity, and that these entity--keyword bindings systematically bias answer selection.  Controlled experiments further demonstrate that such bindings are shaped by frequency imbalance during pretraining.  We also characterize the joint effects of multiple keywords, including both synergistic and competitive interactions.
    \item We propose a keyword-perturbation method for hallucination detection.  By intervening on candidate keywords in embedding space and measuring changes in the likelihood margin between competing answers, the method distinguishes shortcut-supporting cues from constraint-supporting evidence.  It therefore supports both error detection and localization of the evidence driving an error.
    \item We systematically examine hallucination types across multiple benchmarks and characterize which detection signals are appropriate for different failure modes.  Keyword perturbation is most effective for errors driven by competition among local associations, but is less advantageous for failures dominated by distributed multi-step computation.  Conversely, uncertainty- and representation-based detectors are poorly matched to systematic errors driven by misleading keyword associations.
\end{enumerate}

\iffalse
\subsection{Common Experimental Setup}
\label{sec:common-setup}
Unless stated otherwise, the model under study is
\texttt{NousResearch/Meta-Llama-3.1-8B-Instruct}, evaluated in bfloat16 with
its native chat template. We use four frozen candidate-conditioned populations:
1,084 ScientistQA name questions, 1,000 class-balanced TriviaQA questions, 942
class-balanced GSM8K questions, and 1,000 class-balanced DROP questions.
``Correct'' always refers to the original, unperturbed generation. ScientistQA
contains 337 person groups, and supervised evaluations keep questions about
the same person in one fold; remaining benchmarks use source-item IDs.

We score candidates by length-normalized conditional log probability. Given
the wrong candidate $a_w$ and the correct candidate $a_r$, define
\begin{equation}
m(x)=\ell(a_w\mid x)-\ell(a_r\mid x),\qquad
\ell(a\mid x)=|a|^{-1}\sum_{t=1}^{|a|}\log p(a_t\mid x,a_{<t}).
\label{eq:wrong-right-margin}
\end{equation}
Thus $m(x)>0$ favors the wrong candidate.  On erroneous items the original
generation is $a_w$; on correct items it is $a_r$.  We retain this fixed
wrong-minus-right orientation throughout.
\fi

\section{Entity--Keyword Binding Effects}
\label{sec:binding}
Here's a motivating example. Consider a prompt asking the model
to choose between Enrico Fermi and Albert Einstein: “A physicist and university teacher made major contributions to modern physics, but did not formulate the theory of special relativity.” Although
the explicit constraint “did not formulate special
relativity” rules out Einstein and supports Fermi,
the model may still answer “Albert Einstein”.

\subsection{Latent Key Model}
\paragraph{Theory.}
Our perturbation analysis adopts a key-only specialization of the Bayesian
latent-variable setting in \citet{hu2026reasoningpriors}.  Let $k$ denote a
salient keyword activated by the prompt, with all other latent structure
marginalized into the key-conditional prediction $P(a\mid x,k)$.  The resulting
model prediction is
\begin{equation}
P(a\mid x)=\sum_k P(k\mid x)P(a\mid x,k).
\label{eq:latent-key-mixture}
\end{equation}
For a correct key $k_r$ and a shortcut key $k_w$, Bayes' rule decomposes their
posterior odds into a learned prior and prompt-specific evidence:
\begin{equation}
\frac{P(k_w\mid x)}{P(k_r\mid x)}
=
\frac{P(k_w)}{P(k_r)}
\frac{P(x\mid k_w)}{P(x\mid k_r)}.
\label{eq:key-posterior-odds}
\end{equation}
The first ratio reflects associations acquired during pretraining, whereas the
second reflects how the current prompt supports the two keys.  As shown by
\citet{hu2026reasoningpriors}, a sufficiently strong shortcut prior can dominate
even when the prompt contains information supporting the correct answer.

Keyword perturbation tests which part of the prompt sustains these posterior
odds.  Neutralizing a span $s$ changes $x$ to $x\setminus s$ while leaving the
learned prior fixed.  The corresponding change in latent log odds is
\begin{equation}
\Delta_s=
\log\frac{P(k_w\mid x)}{P(k_r\mid x)}
-
\log\frac{P(k_w\mid x\setminus s)}{P(k_r\mid x\setminus s)}.
\label{eq:key-odds-perturbation}
\end{equation}
Because the keys are latent, we estimate this change with the observable
wrong--right likelihood-margin effect $u_s=m(x)-m(x\setminus s)$.  Under the
two-key separation inherited from \citet{hu2026reasoningpriors}, and when the
key-conditional answer distributions vary locally less than the key weights,
$u_s$ is a surrogate for $\Delta_s$.

This view gives simple expected signatures.  For an erroneous response driven
by a shortcut, the wrong key has larger posterior weight and typically
$m(x)>0$; perturbing a keyword that supports this key should reduce the
wrong-answer preference, giving $u_s>0$.  For a correct response, the
constraint-sensitive key dominates and typically $m(x)<0$; perturbing a
keyword that supports the correct key should move the margin toward the wrong
answer, giving $u_s<0$.  Spans that provide little evidence for either key
should have $u_s\approx0$.  These are probabilistic signatures rather than
deterministic rules: an error may depend on several distributed cues, and a
perturbation may also change the key-conditional prediction in
Equation~\ref{eq:latent-key-mixture}.

\paragraph{Experiments.}
We match question phrases to five profile attributes: award received,
education, field, occupation, and position held. For span $s$, the intervention
replaces its input embeddings by the model-wide mean embedding, leaving token
positions and all other tokens unchanged, and measures
$u_s=m(x)-m(x\setminus s)$. A nearby same-width non-attribute span controls for
intervention length and position. The analysis contains 1,069 parse-valid
questions with an eligible attribute (455 erroneous and 614 correct). Effects
are averaged within question; intervals use 10,000 person-group bootstrap draws.

Owner-aligned frequency is a separate free-generation assay over the same
frozen attributes. For each keyword, four attribute-specific paraphrases ask
about its original owner and the other scientist, with five continuations per
paraphrase (temperature 0.8, maximum 32 new tokens, seed 42). Exact $\Delta f$
is the exact-mention-frequency difference, averaged over prompts and keywords
within item and then over items. The assay contains 2,217 item--attribute
records and uses the same 455/614 split and person bootstrap. See Table~\ref{tab:keyword_perturbation_binding} for results.

\subsection{Keyword Binding Effects}
\paragraph{Theory.}
The first factor in Equation~\ref{eq:key-odds-perturbation} is the prompt-specific likelihood ratio, which can be dominate by an entity-keyword binding acquired during pretraining. 
\paragraph{Experiments.}
See Table~\ref{tab:keyword_perturbation_binding} for results.

\begin{table*}[t]
\centering
\caption{Keyword perturbation effects and owner-aligned frequency differences for erroneous and correct responses.}
\label{tab:keyword_perturbation_binding}
\label{tab:free_generation_binding}
\small
\setlength{\tabcolsep}{1.5pt}
\begin{tabular}{lccc|ccc}
\toprule
& \multicolumn{3}{c}{\textbf{Erroneous responses}}
& \multicolumn{3}{c}{\textbf{Correct responses}} \\
\cmidrule(lr){2-4}\cmidrule(lr){5-7}
Attribute
& $n$ & Perturbation & Exact $\Delta f$
& $n$ & Perturbation & Exact $\Delta f$ \\
\midrule
\textbf{Overall}
& \textbf{455}
& $\mathbf{0.451}$\,{\scriptsize $[0.386,\,0.520]$}
& $\mathbf{0.335}$\,{\scriptsize $[0.297,\,0.374]$}
& \textbf{614}
& $\mathbf{0.499}$\,{\scriptsize $[0.443,\,0.561]$}
& $\mathbf{0.319}$\,{\scriptsize $[0.288,\,0.352]$} \\
Award received
& 327
& $0.428$\,{\scriptsize $[0.349,\,0.506]$}
& $0.484$\,{\scriptsize $[0.424,\,0.545]$}
& 377
& $0.439$\,{\scriptsize $[0.365,\,0.516]$}
& $0.445$\,{\scriptsize $[0.390,\,0.501]$} \\
Education
& 144
& $0.648$\,{\scriptsize $[0.542,\,0.758]$}
& $0.287$\,{\scriptsize $[0.231,\,0.342]$}
& 222
& $0.725$\,{\scriptsize $[0.635,\,0.815]$}
& $0.336$\,{\scriptsize $[0.291,\,0.382]$} \\
Field
& 193
& $0.137$\,{\scriptsize $[0.062,\,0.214]$}
& $0.192$\,{\scriptsize $[0.139,\,0.246]$}
& 264
& $0.230$\,{\scriptsize $[0.146,\,0.323]$}
& $0.210$\,{\scriptsize $[0.163,\,0.256]$} \\
Occupation
& 80
& $0.101$\,{\scriptsize $[-0.028,\,0.215]$}
& $0.081$\,{\scriptsize $[0.010,\,0.154]$}
& 149
& $0.213$\,{\scriptsize $[0.124,\,0.307]$}
& $0.146$\,{\scriptsize $[0.087,\,0.208]$} \\
Position held
& 64
& $0.540$\,{\scriptsize $[0.369,\,0.733]$}
& $0.207$\,{\scriptsize $[0.145,\,0.273]$}
& 97
& $0.633$\,{\scriptsize $[0.486,\,0.794]$}
& $0.211$\,{\scriptsize $[0.153,\,0.276]$} \\
\bottomrule
\end{tabular}

\vspace{2pt}
\begin{minipage}{0.98\linewidth}
\footnotesize
\textit{Notes.}
Perturbation denotes the change in the wrong-versus-right answer margin after neutralizing the attribute phrase.
Exact $\Delta f$ is
$P(k\mid\text{original owner})-P(k\mid\text{other person})$
under exact keyword matching.
Bracketed values are person-clustered bootstrap 95\% confidence intervals.
\end{minipage}
\end{table*}

\subsection{Binding Effects Are Shaped by Pretraining Frequency}
\paragraph{Theory.}
The second factor in Equation~\ref{eq:key-odds-perturbation} is the prior odds of the competing keys, which can be shaped by pretraining frequency.  We test this hypothesis with a controlled frequency--dose experiment. 
\paragraph{Experiments.}
See Table~\ref{tab:binding-dose-three-models} for results.

We use 50 mirrored pairs of fictional people and facts $B$/$F$. For one person,
$B$ occurs $0.05N$ times and $F$ occurs $(0.05+0.65d)N$ times; assignment is
reversed for the other, with $N=40$ and $d\in\{0,.25,.5,.75,1\}$. Remaining
examples contain a neutral fact; forward/reverse templates alternate by fact.

At each dose we reload the checkpoint and fine-tune only its last four blocks
and final normalization layer for two epochs. We use AdamW (learning rate
$10^{-4}$, zero weight decay), batch size 12, bfloat16, and gradient clipping
at 1.0. Checkpoints are Qwen2.5-3B-Instruct, Llama-3.2-3B-Instruct, and
Ministral-3-3B-Instruct-2512. Evaluation averages both profile orders and
computes $\Delta_{\mathrm{B-F}}=m(x_B)-m(x_F)$ from full-name likelihoods. The
slope is OLS over five doses. The table gives means and runs for seeds 42--44.

\begin{table*}[t]
\centering
\caption{Frequency--dose response across models.  Values are candidate-preference differences $\Delta_{\mathrm{B-F}}$; larger values indicate stronger transfer of asymmetric exposure frequency into candidate preference.}
\label{tab:binding-dose-three-models}
\small
\setlength{\tabcolsep}{1.5pt}
\newcommand{\seedcell}[4]{%
  \begin{tabular}[c]{@{}r@{\hspace{2pt}}r@{}}
  & {\scriptsize #2} \\
  \textbf{#1} & {\scriptsize #3} \\
  & {\scriptsize #4}
  \end{tabular}}
\begin{tabular}{lcccccc}
\toprule
Model & $d=0$ & $d=.25$ & $d=.5$ & $d=.75$ & $d=1$ & Slope \\
\midrule
Qwen2.5-3B
& \seedcell{.028}{.022}{.024}{.039}
& \seedcell{.055}{.028}{.050}{.086}
& \seedcell{.131}{.171}{.073}{.149}
& \seedcell{.200}{.346}{.096}{.156}
& \seedcell{.232}{.286}{.199}{.210}
& \seedcell{.221}{.338}{.159}{.165} \\
\midrule
Llama-3.2-3B
& \seedcell{.002}{-.005}{.000}{.011}
& \seedcell{.013}{.008}{.029}{.003}
& \seedcell{.014}{.009}{.016}{.018}
& \seedcell{.024}{.009}{.029}{.033}
& \seedcell{.064}{.047}{.124}{.021}
& \seedcell{.054}{.042}{.099}{.020} \\
\midrule
Ministral-3-3B
& \seedcell{.013}{-.031}{.020}{.050}
& \seedcell{.059}{.074}{.062}{.042}
& \seedcell{.095}{.143}{.121}{.021}
& \seedcell{.285}{.374}{.202}{.279}
& \seedcell{.564}{.589}{.398}{.704}
& \seedcell{.531}{.616}{.359}{.618} \\
\bottomrule
\end{tabular}

\vspace{2pt}
\begin{minipage}{0.96\linewidth}
\footnotesize
\textit{Notes.}
Bold values are means across three random seeds.  Smaller values to their right report seeds 42, 43, and 44 from top to bottom.
\end{minipage}
\end{table*}

\subsection{Keywords associations}
\label{sec:multi}
\paragraph{Theory.}
The latent key model in Section~\ref{sec:binding} assumes that a single keyword dominates the posterior odds.  In practice, multiple keywords may be present and interact in complex ways. 
\paragraph{Experiments.}
We therefore examine the joint effects of pairs of keywords on the wrong--right margin $m(x)$. See Table~\ref{tab:pairwise-keyword-interactions} and Table~\ref{tab:joint-keyword-interventions} for results.

We enumerate the complete perturbation factorial for all 1,084 ScientistQA name
questions. Eight invalid-candidate questions and 16 with fewer than two usable
keywords are excluded, leaving 1,060 questions (611 correct, 449 erroneous) and
4,128 pairs. The local interaction is the second finite difference of $m$; its
Banzhaf counterpart averages this difference over all backgrounds. At the
prespecified $\tau=0.10$, synergy/redundancy require local and Banzhaf effects
to agree in sign and exceed $\tau$; opposite-sign single effects with minimum
magnitude above $\tau$ define competition. Higher-order mass is the fraction
of absolute Harsanyi mass at order three or above. Rates are question-weighted;
intervals use 10,000 question bootstrap draws.

\begin{table*}[t]
\centering
\caption{Distribution of pairwise keyword interactions in correct and erroneous responses.}
\label{tab:pairwise-keyword-interactions}
\small
\setlength{\tabcolsep}{5pt}
\begin{tabular}{lrrccccc}
\toprule
Subset & $n_q$ & $n_{\mathrm{pair}}$ & Competition & Redundancy & Synergy & Mixed & Higher-order \\
\midrule
Correct responses
& 611 & 2,347 & 31.2\% & 16.2\% & 3.0\% & 49.7\% & 28.5\% \\
Erroneous responses
& 449 & 1,781 & 37.3\% & 16.1\% & 2.1\% & 44.5\% & 29.9\% \\
\bottomrule
\end{tabular}

\vspace{2pt}
\begin{minipage}{0.96\linewidth}
\footnotesize
\textit{Notes.}
$n_q$ is the number of questions and $n_{\mathrm{pair}}$ is the number of evaluated keyword pairs.  Interaction percentages are computed over keyword pairs; minor deviations from 100\% are due to rounding.
\end{minipage}
\end{table*}

\begin{table}[t]
\centering
\caption{Answer probabilities under targeted keyword removal and matched-random controls.}
\label{tab:joint-keyword-interventions}
\small
\setlength{\tabcolsep}{6pt}
\begin{tabular}{lcc}
\toprule
Intervention & $P(\mathrm{right})$ & $P(\mathrm{wrong})$ \\
\midrule
None                                  & 37.9\% & 60.9\% \\
Remove wrong-answer keyword           & 52.4\% & 44.7\% \\
Remove correct-answer keyword         & 15.4\% & 83.5\% \\
Remove both keywords                  & 25.9\% & 70.5\% \\
\midrule
Matched random (wrong-answer side)    & 38.0\% & 60.1\% \\
Matched random (correct-answer side)  & 32.5\% & 65.9\% \\
Matched random (both sides)           & 27.7\% & 70.4\% \\
\bottomrule
\end{tabular}

\vspace{2pt}
\begin{minipage}{0.94\linewidth}
\footnotesize
\textit{Notes.}
This frozen confirmation contains 91 baseline errors. Person groups are reserved
by SHA-256 parity before outcomes are inspected; an opposite-sign pair with
minimum one-sided effect above 0.10 is selected by strongest minimum effect.
Probabilities are exact-name frequencies over ten continuations (temperature
0.8, at most 20 new tokens) under mean-embedding neutralization. All conditions
share a per-question seed; placebo spans are fixed in advance and matched in
width and question position.
\end{minipage}
\end{table}




\section{Hallucination Detection via Keyword Perturbation}
\label{sec:detection}
\subsection{Mechanism}
\paragraph{Probe and distinguishability.}
A question can contain both a span that supports the shortcut key $k_w$ and a
span that constrains the model toward the correct key $k_r$, regardless of
whether the final response is correct. Their presence, and even the magnitude
of their individual effects, is therefore not sufficient to determine
correctness. What matters is the balance of evidence that they provide. If a
span supports $k_w$, neutralizing it reduces the posterior preference for the
shortcut key and moves the wrong--right margin toward the correct answer,
yielding $u_s>0$. Conversely, neutralizing a span that supports $k_r$ removes
evidence for the correct association and moves the margin toward the wrong
answer, yielding $u_s<0$. A correct response occurs when the constraint-based
evidence is strong enough to overcome the competing shortcut preference; an
erroneous response occurs when it is not. The joint pattern of signed
perturbation effects thus reveals the residual preference that remains after
the relevant evidence is removed.

This interpretation requires changes in the latent-key posterior to produce
observable changes in the answer margin. For a probe-known question, the model
can recover the queried attribute when conditioned on an appropriate key: in
Equation~\ref{eq:latent-key-mixture}, $P(a_r\mid x,k_r)$ is concentrated on the
correct answer and is distinguishable from the answer distribution induced by
$k_w$. The perturbation effect $u_s$ can then serve as a reliable readout of
which association a span supports.

For a probe-unknown question, even the correct key does not anchor a
concentrated distribution over the correct answer. The key-conditional answer
distributions consequently overlap, so a perturbation may change $P(k\mid x)$
without producing a consistent change in $m(x)$. Correct-supporting and
shortcut-supporting spans then become difficult to distinguish from their
perturbation effects: not because they represent the same key, but because the
mapping from key selection to answer preference is weak. We therefore restrict
the present detection analysis to probe-known questions, for which the signed
effects have a well-defined mechanistic interpretation.

\paragraph{Detection setup.}
For each candidate keyword, we neutralize its embedding and measure the
resulting change in the wrong--right margin $m(x)$. We use the resulting vector
of perturbation effects to train a simple classifier that distinguishes correct
from erroneous responses. We consider three strategies for selecting candidate
keywords: exact selection enumerates all spans, whereas attention and gradient
selection screen spans using candidate-conditioned attention and input
gradients, respectively.

\subsection{Empirical Evaluation}
The detector summarizes likelihood and representation responses to mean-embedding
neutralization. Exact selection enumerates every span; attention and gradient
selection use candidate-conditioned attention or input gradients only to screen
spans. The 127-dimensional vector contains 47 scalar features, four 8-dimensional
PCA blocks from base and effect-weighted states for both candidates, and a
48-dimensional PCA projection of the layer-14 mean state. Standardization and
PCA are fold-local. A balanced logistic regression ($C=0.03$, \texttt{liblinear},
5,000 iterations) gives the error score.

In-domain AUROCs average seeds 42--44 under five-fold OOF prediction, with
stratified group folds except for TriviaQA unique keys. No benchmark-specific
tuning is performed. Transfer fits all preprocessing and the classifier on
1,084 ScientistQA items, freezes them, and applies them to 477 athlete, building,
and musician questions. See Table~\ref{tab:perturbation-detection-summary} for results.

\begin{table*}[t]
\centering
\caption{AUROC of keyword-perturbation detection across models and benchmarks.}
\label{tab:perturbation-detection-summary}
\small
\setlength{\tabcolsep}{6pt}
\begin{tabular}{llccccc}
\toprule
Model & Selection & ScientistQA$^\star$ & TriviaQA & GSM8K & DROP & Transfer \\
\midrule
 & Exact     & .910 & .949 & .7473 & .919 & .931 \\
Llama & Attention & .906 & .946 & .7430 & .915 & .922 \\
 & Gradient  & .905 & .943 & .7426 & .915 & .928 \\
\midrule
 & Exact     & .822 & .964 & .7439 & .968 & .850 \\
Mistral & Attention & .817 & .964 & .7372 & .969 & .803 \\
 & Gradient  & .819 & .966 & .7415 & .965 & .809 \\
\midrule
 & Exact     & .861 & .909 & .7603 & .907 & .848 \\
Qwen & Attention & .835 & .902 & .7575 & .902 & .761 \\
 & Gradient  & .830 & .898 & .7599 & .884 & .811 \\
\bottomrule
\end{tabular}

\vspace{2pt}
\begin{minipage}{0.92\linewidth}
\footnotesize
\textit{Notes.}
In-domain results use repeated five-fold out-of-fold evaluation; Multidomain reports aggregate transfer performance from a detector frozen on ScientistQA.  
$\star$ means using the probe-known subset of ScientistQA. Full results are provided in Appendix~\ref{app:detection}.
\end{minipage}
\end{table*}

We tested four different families of hallucination detection methods, see Table~\ref{tab:detection-auroc} for results. 

\begin{table*}[t]
\centering
\caption{Hallucination-detection AUROC across method families and benchmarks.}
\label{tab:detection-auroc}
\small
\setlength{\tabcolsep}{4pt}
\begin{tabular}{llccccc}
\toprule
Family & Method & ScientistQA$^\star$ & Full ScientistQA & TriviaQA & GSM8K & DROP \\
\midrule
& Ours (Exact) & \textbf{.910} & .794 & \textbf{.949} & .7473 & \textbf{.919} \\
Perturbation
& Ours (Attention) & .906 & -- & .946 & .7430 & .915 \\
& Ours (Gradient) & .905 & -- & .943 & .7426 & .915 \\
\midrule
Representation
& Aiersilan-style & .904 & \textbf{.810} & .879 & .744 & .878 \\
& SAPLMA & .671 & .657 & .840 & .719 & .869 \\
& ICR Probe & .573 & .569 & .726 & .707 & .837 \\
\midrule
Uncertainty
& SelfCheckGPT-NLI & .593 & .580 & .747 & \textbf{.790} & .765 \\
& SeSE & .560 & -- & .705 & .632 & .749 \\
& Semantic Entropy & .518 & .528 & .708 & .503 & .752 \\
\midrule
Evidence
& MiniCheck (unilateral) & .717$^\dagger$ & .809$^\dagger$ & .658 & .558 & .604 \\
& MiniCheck (contrastive) & .926$^\dagger$ & .987$^\dagger$ & .823 & .506 & .526 \\
\bottomrule
\end{tabular}

\vspace{2pt}
\begin{minipage}{0.9\linewidth}
\footnotesize
\textit{Notes.}
Bold values indicate the best result on each benchmark. For Full ScientistQA,
MiniCheck obtains AUPRCs of .749 (unilateral) and .985 (contrastive); the
contrastive variant achieves .957 balanced accuracy at a fixed threshold.
$\dagger$: provided with full profiles of the scientists after answer generation.
\end{minipage}
\end{table*}

\iffalse
\paragraph{Standard detector baselines.}
All methods use the same frozen generations and labels. The Aiersilan-style
probe uses answer-final and answer-mean layer-14 states, separate fold-local
PCA-8 projections, and balanced logistic regression. SAPLMA uses the layer-28
final-answer state and the published 4096--256--128--64--1 ReLU MLP (Adam, BCE,
five epochs, batch 32). ICR uses layerwise Jensen--Shannon divergences from the
top-20 attention-selected tokens and its published 128/64/32 network (batch
normalization, Leaky-ReLU 0.01, dropout 0.3, Adam at $5\times10^{-4}$, up to
50 epochs with fold-internal validation). Supervised probes use three seeds and
five group-disjoint folds.

The $K=6$ method draws six samples at temperature 0.7 and top-$p=0.95$ and
scores disagreement with the canonicalized greedy answer. Semantic Entropy
uses the paper protocol: ten samples at temperature 1.0, top-$p=0.9$, top-$k=50$;
bidirectional DeBERTa-large-MNLI entailment defines semantic classes, and
entropy is computed over summed class probability masses. MiniCheck uses
MiniCheck-Flan-T5-Large. TriviaQA/DROP become declarative claims; GSM8K is
checked sentence-wise using minimum support. Unilateral is one minus generated
support; contrastive is alternative minus generated support and can therefore
use a gold-derived manifest candidate.
\fi

\section{No Universe Winner}
\label{sec:category}
In this section, we systematically examine why different detection methods succeed or fail on different benchmarks.  We find that the effectiveness of a given method depends on the type of error, which may help explain why detectors do not always transfer reliably across benchmarks. 

\subsection{Category Definitions}
We use four complementary signals---uncertainty (U), external-evidence support (E), keyword perturbation (P), and internal representations (R)---to define an operational taxonomy of erroneous responses.  These categories describe observable diagnostic patterns rather than uniquely identified causal mechanisms.

All axes are joined by item key to one baseline generation. Within benchmark,
U-high/U-low and P-high/E-high use inclusive upper/lower 30th-percentile value
cutoffs over the complete population; R is correct-like when its group-OOF error
probability is below 0.5. Evidence-completion labels exist for ScientistQA Names
and TriviaQA; Profiles uses grounded support, while GSM8K and DROP have no E
label. DROP context dependence lacks an independent matched-placebo confirmation.
These are multi-label operational signatures, not human causal annotations.
For GSM8K uncertainty, the later full run uses the original natural chain-of-
thought prompt on all 942 items, draws six continuations at temperature 0.7 and
top-$p=0.95$ (maximum 192 new tokens), computes canonical-numeric entropy from
the first three samples, and reserves the last three for majority-answer
validation. It reuses the identically configured 300-item pilot and samples the
remaining 642 items with seed 20260823.

\begin{description}[leftmargin=0pt,labelindent=0pt,itemsep=2pt,topsep=3pt]
\item[Knowledge deficit.]
$U$-high $\wedge$ $E$-high $\wedge$ evidence-completion gain $>0$.
The answer is unstable and poorly supported, while supplying relevant evidence increases preference for the correct answer.

\item[Confident unsupported error.]
$\neg U$-high $\wedge$ $E$-high $\wedge$ evidence-completion gain $>0$.
The answer lacks evidential support without showing substantial uncertainty, but improves when relevant evidence is supplied.

\item[Context-induced local dependence.]
$P$-high $\wedge$ targeted intervention outperforms a matched placebo.
Perturbing a small set of relevant input phrases weakens the wrong answer more than perturbing matched irrelevant phrases.

\item[Reasoning instability.]
$U$-high.
Repeated sampling on the same input produces high answer entropy.  This indicates an unstable decision, but does not by itself localize a faulty reasoning step.

\item[Stable self-consistent error.]
$U$-low $\wedge$ $R$ correct-like.
The model repeatedly produces the same wrong answer even though an out-of-fold representation scorer assigns its hidden-state trajectory a correct-like score.

\item[Mixed.]
The response satisfies at least two of the criteria above.

\item[Unclassified.]
The response satisfies none of the current criteria.  This indicates a limitation of the present UEPR coverage, not the absence of an underlying mechanism.
\end{description}

\begin{table*}[t]
\centering
\caption{Distribution of operational hallucination types under unified UEPR scoring.}
\label{tab:unified-hallucination-types}
\small
\setlength{\tabcolsep}{3pt}
\begin{tabular}{lccccc}
\toprule
Hallucination type & ScientistQA & ScientistQA$^\dagger$ & TriviaQA & GSM8K & DROP \\
\midrule
Knowledge deficit
& 93 (20.5\%) & -- & 50 (10.0\%) & -- & -- \\
Confident unsupported error
& 146 (32.2\%) & 34 (27.9\%) & 25 (5.0\%) & -- & -- \\
Context-induced dependence
& 38 (8.4\%) & 7 (5.7\%) & 217 (43.4\%) & 63 (13.4\%) & 275 (55.0\%) \\
Reasoning instability
& 159 (35.1\%) & 83 (68.0\%) & 310 (62.0\%) & 248 (52.7\%) & 348 (69.6\%) \\
Stable self-consistent error
& 29 (6.4\%) & 16 (13.1\%) & 85 (17.0\%) & 92 (19.5\%) & 51 (10.2\%) \\
Mixed
& 121 (26.7\%) & 19 (15.6\%) & 220 (44.0\%) & 47 (10.0\%) & 226 (45.2\%) \\
Unclassified
& 122 (26.9\%) & 3 (2.5\%) & 40 (8.0\%) & 115 (24.4\%) & 52 (10.4\%) \\
\bottomrule
\end{tabular}

\vspace{2pt}
\begin{minipage}{0.98\linewidth}
\footnotesize
\textit{Notes.}
$\dagger$: provided with full profiles of the scientists.  The numbers of erroneous responses are 453 for Names, 122 for Profiles, 500 for TriviaQA, 471 for GSM8K, and 500 for DROP.
Categories are multi-label and therefore do not sum to 100\%.
\end{minipage}
\end{table*}

\subsection{Methods Ensemble}
The previous section defines a taxonomy of hallucination types based on the four complementary signals, each can be detected by a different group families.
\begin{table*}[t]
\centering
\caption{Correspondence between failure modes, hallucination types, and diagnostic signals.}
\label{tab:failure-mode-diagnostic-signals}
\small
\setlength{\tabcolsep}{8pt}
\begin{tabular}{lll}
\toprule
Failure locus & Hallucination type & Diagnostic signal \\
\midrule
Incorrect $P(a\mid x,k_r)$ & Knowledge deficit & $E$ \\
Excessive prior odds & Stable self-consistent error & $P$ (weak); not $R$ or $U$ \\
Misread evidence term & Context-induced error & $P$ \\
Diffuse posterior over $k$ & Reasoning instability & $U$ \\
\bottomrule
\end{tabular}
\end{table*}

In this section we turn our attention to the entire ScientistQA benchmark, rather than the probe-known subset. We find that the four families of detection methods are complementary, and that their effectiveness depends on the type of error. 

\begin{table*}[t]
\centering
\caption{Detection performance of perturbation, representation, and fused signals across benchmarks.}
\label{tab:cross-benchmark-signal-fusion}
\label{tab:scientistqa-signal-fusion}
\small
\setlength{\tabcolsep}{9pt}
\begin{tabular}{llccc}
\toprule
Benchmark & Method & AUROC & AUPRC & Balanced accuracy \\
\midrule
& P   & .794 & .773 & .707 \\
ScientistQA & R   & .810 & .832 & .723 \\
& P+R & \textbf{.850} & \textbf{.832} & \textbf{.764} \\
\midrule
& P   & \textbf{.976} & .967 & \textbf{.925} \\
TriviaQA & R   & .893 & .897 & .808 \\
& P+R & .975 & \textbf{.968} & .917 \\
\midrule
& P   & .752 & .767 & \textbf{.702} \\
GSM8K & R   & .733 & .760 & .656 \\
& P+R & \textbf{.765} & \textbf{.784} & .697 \\
\midrule
& P   & .914 & .928 & .828 \\
DROP & R   & .866 & .903 & .815 \\
& P+R & \textbf{.924} & \textbf{.937} & \textbf{.842} \\
\bottomrule
\end{tabular}
\end{table*}

\section{Conclusion}
In this work, we have investigated the gap between model knowledge and its ability to use that knowledge correctly, particularly in the context of hallucinations. We introduced a keyword perturbation method that allows us to identify which keywords contribute to hallucinations and which support correct answers. Our experiments across multiple models and benchmarks demonstrate that this method is effective for detecting hallucinations and attributing errors to specific keywords.
